\documentclass[11pt,a4paper]{article}
\usepackage[margin=0.5in]{geometry}
\usepackage[T1]{fontenc}
\usepackage[utf8]{inputenc}
\usepackage{xcolor}
\usepackage{enumitem}
\usepackage{titlesec}
\usepackage{hyperref}
\usepackage{mdframed}

\ifdefined\pdfgentounicode
  \pdfgentounicode=1
\fi

\definecolor{dark}{RGB}{51, 51, 51}
\definecolor{rulecolor}{RGB}{180, 180, 180}
\definecolor{subbg}{RGB}{245, 245, 245}

% ── Section headings ────────────────────────────────────────────────────────
\titleformat{\section}
    {\large\bfseries\color{dark}}
    {}{0em}{}[\titlerule]
\titlespacing{\section}{0pt}{14pt}{7pt}

\definecolor{linkblue}{RGB}{0, 80, 180}

\hypersetup{
    colorlinks=true,
    linkcolor=linkblue,
    urlcolor=linkblue,
    pdfauthor={Kritarth Dandapat},
    pdftitle={Kritarth Dandapat CV}
}

\setlength{\parindent}{0pt}
\setlength{\parskip}{0pt}

% ── List style ──────────────────────────────────────────────────────────────
\newlist{cvlist}{itemize}{2}
\setlist[cvlist,1]{label=\textbullet, nosep, leftmargin=1.4em,
    itemsep=3pt, topsep=3pt, parsep=0pt, partopsep=0pt}
\setlist[cvlist,2]{label=\textbullet, nosep, leftmargin=1.4em,
    itemsep=3pt, topsep=3pt, parsep=0pt, partopsep=0pt}

% ── Thin divider between entries inside a section ───────────────────────────
\newcommand{\entrydivider}{%
  \vspace{4pt}%
  {\color{rulecolor}\hrule height 0.3pt}%
  \vspace{4pt}%
}

% ── Role / org header ───────────────────────────────────────────────────────
\newcommand{\roleentry}[4]{%
  \vspace{0.5em}%
  \begin{tabular*}{\textwidth}{@{\extracolsep{\fill}} ll}
    \textbf{#1} & #2 \\
    \textit{#3} & \textit{#4} \\
  \end{tabular*}%
  \vspace{0.15em}%
}

% ── Subproject header (lightly shaded pill) ──────────────────────────────────
\newcommand{\subrole}[2]{%
  \vspace{0.45em}%
  \noindent\colorbox{subbg}{%
    \begin{minipage}{\dimexpr\textwidth-2\fboxsep}%
      \begin{tabular*}{\linewidth}{@{\extracolsep{\fill}} ll}
        \textbf{#1} & \textit{\small #2}%
      \end{tabular*}%
    \end{minipage}%
  }%
  \vspace{0.15em}%
}

% ── Project header ───────────────────────────────────────────────────────────
\newcommand{\projectentry}[2]{%
  \vspace{0.55em}%
  \begin{tabular*}{\textwidth}{@{\extracolsep{\fill}} ll}
    \textbf{#1} & \textit{#2} \\
  \end{tabular*}%
  \vspace{0.1em}%
}

% ── Scope sentence ───────────────────────────────────────────────────────────
\newcommand{\scopesentence}[1]{%
  \vspace{0.15em}%
  {\itshape #1}\par
  \vspace{0.1em}%
}

% ── Tech stack tag line ──────────────────────────────────────────────────────
\newcommand{\techstack}[1]{%
  \vspace{2pt}%
  \noindent{\small\textcolor{dark!70}{\textbf{Stack:} #1}}\par
}

\begin{document}

% ════════════════════════════════════════════════════════════════════════════
%  HEADER
% ════════════════════════════════════════════════════════════════════════════
\begin{center}
  {\LARGE\bfseries\color{dark} Kritarth Dandapat}\\[0.3em]
  {\small Incoming PhD Student, Michigan State University \textbar\
         Undergraduate Researcher, University at Buffalo}\\[0.2em]
  {\small \href{mailto:contact@dkritarth.com}{contact@dkritarth.com}
   \textbar\ +1 (716) 612-0016
   \textbar\ Buffalo, NY}\\[0.2em]
  {\small
   \href{https://dkritarth.com}{dkritarth.com} \textbar\
   \href{https://github.com/dkritarth}{github.com/dkritarth} \textbar\
   \href{https://www.linkedin.com/in/kritarth-dandapat}{linkedin.com/in/kritarth-dandapat}}
\end{center}
\vspace{0.3em}

% ════════════════════════════════════════════════════════════════════════════
%  PROFESSIONAL SUMMARY
% ════════════════════════════════════════════════════════════════════════════
\section*{Professional Summary}
Incoming PhD student in Computer Science at Michigan State University (Fall 2026), joining the Data Mining Laboratory under Dr.\ Pang-Ning Tan to work on spatiotemporal machine learning, deep learning-based weather forecasting, and AI adversarial robustness. Currently completing an accelerated three-year BS in Computer Science (Minor: Statistics) at the University at Buffalo, 3.8+ GPA.

Combines research and engineering: co-authored work in digital health and computational materials science, and has independently built and deployed production ML systems spanning mobile health, orthodontic monitoring, and telerehabilitation. Recognized with the PEARL undergraduate research award (\$2,500) and second place at the UB Health Futures Challenge (\$1,000).

% ════════════════════════════════════════════════════════════════════════════
%  EDUCATION
% ════════════════════════════════════════════════════════════════════════════
\section*{Education}

\roleentry{PhD in Computer Science}{Starting Fall 2026}{Michigan State University}{East Lansing, MI}
\begin{cvlist}
  \item Department of Computer Science and Engineering; Data Mining Laboratory (Advisor: \href{https://www.cse.msu.edu/~ptan/}{Dr.\ Pang-Ning Tan})
  \item Planned research focus: spatiotemporal machine learning, deep learning-based weather forecasting (DLWF), and AI adversarial robustness
\end{cvlist}

\entrydivider

\roleentry{Bachelor of Science in Computer Science, Minor in Statistics}{August 2023 -- June 2026}{University at Buffalo, SUNY}{Buffalo, NY}
\begin{cvlist}
  \item Specialization in artificial intelligence; coursework spanning reinforcement learning, computer vision, machine learning, quantum computing, and distributed systems
  \item GPA 3.8+; Dean's List all semesters; accelerated three-year completion (21--22 credit semesters)
  \item Presidential Scholarship: \$15,000 per year; PEARL Award: \$2,500 competitive undergraduate research grant (Experiential Learning Network)
\end{cvlist}

% ════════════════════════════════════════════════════════════════════════════
%  RESEARCH EXPERIENCE
% ════════════════════════════════════════════════════════════════════════════
\section*{Research Experience}

\roleentry{Research Assistant}{June 2024 -- Present}{Embedded Sensing and Computing (ESC) Group, University at Buffalo}{Buffalo, NY}
\scopesentence{Mobile sensing for healthcare and rehabilitation with \href{https://cse.buffalo.edu/~wenyaoxu/index.html}{Prof.\ Wenyao Xu}.}

\subrole{OralScan: AI-Assisted Geriatric Oral Screening}{June 2024 -- February 2025}
\begin{cvlist}
  \item Architected a full-stack mobile health platform (React Native, React web) integrating end-to-end \textbf{YOLOv8} vision pipelines for guided intraoral image acquisition, multi-class dental disease screening, and automated tooth numbering across temporal scan sessions.
  \item Formulated a spatio-temporal scan-guidance system that aggregates sequential intraoral frames to produce consistent disease-detection signals robust to user-induced motion variability.
  \item Co-authored formative usability and acceptability study submitted to \textit{Smart Health} (under review); delivered oral presentations to clinical stakeholders at CTSI and SURC 2024.
  \item Awarded second place (\$1,000), University at Buffalo Health Futures Challenge (Spring 2026).
\end{cvlist}

\subrole{OrthoScan: Orthodontic Remote Monitoring}{February 2025 -- October 2025}
\begin{cvlist}
  \item Engineered mobile and backend systems for at-home orthodontic progress tracking, deploying \textbf{YOLO}-based hardware detection paired with depth-assisted spatial measurement pipelines to quantify structural change across longitudinal patient visits.
  \item Reconstructed and analyzed \textbf{3D dental structures} from intraoral scans using pre-trained \textbf{GNNs} for 3D mesh processing; computed quantitative tooth position changes in both 2D image space and 3D world coordinates to track orthodontic progress with geometric precision.
  \item Translated sequential intraoral imagery into calibrated geometric progress signals, enabling quantitative remote monitoring between in-office appointments and reducing reliance on subjective clinician assessment.
\end{cvlist}

\subrole{mRehab: Telerehabilitation Platform (\href{https://mrehab.agency/}{mrehab.agency})}{November 2025 -- Present}
\begin{cvlist}
  \item Provisioned and configured \textbf{AWS} backend infrastructure to host testing and staging environments, enabling AI agent capabilities to operate within the live application stack.
  \item Deployed a real-time live sensor logging site for continuous ingest and monitoring of wearable sensor streams during active therapy sessions.
  \item Implemented \textbf{Kalman filters} for sensor fusion and activity detection, improving signal smoothness and robustness of exercise recognition under noisy real-world IMU conditions.
  \item Contributed to session logging, multi-modal sensor ingest pipelines, and exercise scoring algorithms deployed in live outpatient therapy workflows.
\end{cvlist}

\entrydivider

\roleentry{Undergraduate Researcher}{June 2025 -- Present}{Peng Research Lab, University at Buffalo}{Buffalo, NY}
\scopesentence{Computational materials science and clean-energy applications with \href{https://ubwp.buffalo.edu/jiayu-peng-lab/}{Prof.\ Jiayu Peng}.}

\subrole{Perovskite Ordering and Symmetry-Aware Graph Neural Networks}{June 2025 -- November 2025}
\begin{cvlist}
  \item Benchmarked symmetry-aware GNN architectures (including \textbf{ALIGNN}) on ordering-dependent energetics in crystalline perovskites; contributed benchmarking results and revised analysis to arXiv:2409.13851 and the open-source \texttt{PerovskiteOrderingGCNNs} repository.
  \item Migrated experiment tracking infrastructure from SigOpt to \textbf{Weights \& Biases}; automated training loops, hyperparameter sweeps, and multi-architecture evaluation pipelines, reducing manual overhead and improving reproducibility.
  \item Co-authored commentary on agentic AI for multimetallic catalyst discovery, published on \textit{ChemRxiv} (2025).
\end{cvlist}

\subrole{Equivariant MLIPs and Alchemical Extensions}{December 2025 -- Present}
\begin{cvlist}
  \item Extended a dual-topology alchemical graph construction -- originally built for \textbf{MACE} -- to two additional equivariant MLIP architectures (fairchem \textbf{eSEN}, Meta \textbf{UMA}), enabling differentiable composition interpolation over disordered structures without retraining.
  \item Designed a $\lambda$-weighted message-passing and expert-routing scheme for UMA's \textbf{Mixture-of-Linear-Experts} backbone, correcting double-counting in composition-dependent expert gating and adding an endpoint-collapse mechanism for exact pure-composition limits.
  \item Discovered and fixed a reference-energy bug causing non-physical $\sim$15\,eV energy discontinuities near pure compositions -- invisible to force-based relaxation but critical for free-energy calculations -- and added regression tests to prevent recurrence.
  \item Built a 12+ test automated verification suite (energy/force/stress/gradient agreement to $\sim 10^{-7}$\,eV) plus benchmark figures and formula-level documentation cross-validating \textbf{MACE}, \textbf{eSEN}, and \textbf{UMA}.
\end{cvlist}


% ════════════════════════════════════════════════════════════════════════════
%  PUBLICATIONS AND PREPRINTS
% ════════════════════════════════════════════════════════════════════════════
\section*{Publications and Preprints}
\begin{cvlist}
  \item \textbf{Dandapat, K.}, Emily [Surname], Master, T. A., Das, A., Bo, W., Cavuoto, L., and Xu, W. (2026). From Community Feedback to Measurement Redesign: Iterating mRehab for Accessible Home-Based Stroke Rehabilitation. \textit{HumanSys 2026}. Preprint.

  \item \textbf{Dandapat, K.}, Master, T. A., Das, A., Bo, W., Lei, M., Cavuoto, L., Subryan, H., and Xu, W. (2026). Separating Safety from Preference: A Role-Gated Evaluation Platform for Clinical Rehabilitation LLMs. \textit{HumanSys 2026}. Preprint.

  \item Lei, M., Das, A., Master, T. A., \textbf{Dandapat, K.}, Reddipogu, P., Xian, J., Rowe, V., Craft, L., Tabb, K., Cavuoto, L., Bhattacharjya, S., Jo, H. J., Subryan, H., Bo, W., and Xu, W. (2026). Towards AI Agents for Intelligent Voice-Driven Interaction in a Home-Based Stroke Rehabilitation System. University at Buffalo and Georgia State University. Preprint.

  \item \textbf{Dandapat, K.}, et al. (2026). Demo: mRehab -- A Tangible, Offline-First Mobile Platform for Post-Stroke Upper-Limb Rehabilitation with an Edge-Cloud Voice [Assistant]. \textit{MobiComm 2026} (Demo Track). Preprint.

  \item Soni, P., \textbf{Dandapat, K.}, Gherardi, A., Bo, W., Li, R., and Xu, W. (2025). OralScan, an AI-Powered Mobile Tool for Geriatric Oral Healthcare: A Formative Usability and Acceptability Study. Submitted to \textit{Smart Health} (under review).

  \item Peng, J., Liu, C., Luo, Y., and \textbf{Dandapat, K.} (2025). Accelerating Multimetallic Catalyst Discovery with Robotics and Agentic AI. \textit{ChemRxiv}, version 1. DOI: \href{https://doi.org/10.26434/chemrxiv-2025-13n3f}{10.26434/chemrxiv-2025-13n3f}.

  \item Peng, J., et al.\ (2024). Learning Ordering in Crystalline Materials with Symmetry-Aware Graph Neural Networks. \textit{arXiv} preprint \href{https://arxiv.org/abs/2409.13851}{arXiv:2409.13851}. (Contributed benchmarking and revisions.)
\end{cvlist}

% ════════════════════════════════════════════════════════════════════════════
%  RESEARCH PRESENTATIONS
% ════════════════════════════════════════════════════════════════════════════
\section*{Research Presentations}
\begin{cvlist}
  \item Oral presentation and live demo -- CTSI Research Group of Doctors, Nurses, and Medical Students (April 2024)
  \item Oral presentation -- SUNY Undergraduate Research Conference (SURC) (April 2024)
\end{cvlist}

% ════════════════════════════════════════════════════════════════════════════
%  SELECTED PROJECTS
% ════════════════════════════════════════════════════════════════════════════
\section*{Selected Projects}

\projectentry{\href{https://dkritarth.com/PaddockPsychRL/}{PaddockPsychRL} \textnormal{[\href{https://github.com/Kritarth-Dandapat/PaddockPsychRL}{GitHub}]}}{Multi-Agent RL}
\begin{cvlist}
  \item Formulated a multi-agent RL framework integrating Formula 1 temporal performance signals into agent psychology models within \textbf{PettingZoo} cooperative environments, mapping real-world behavioral dynamics to learned policy structure.
  \item Constructed driver-specific agent profiles from \textbf{FastF1} 2025 lap-time consistency metrics, encoding temporal performance variability as per-agent action-noise scales to reflect authentic decision-making under pressure.
  \item Benchmarked tabular Q-learning and \textbf{MAPPO} baselines against psych-conditioned policies; psych-aware agents improved Strategy Resilience Score from 0.54 to 0.63 (+16.7\%) under noisy multi-agent coordination.
\end{cvlist}
\techstack{Python, FastF1, PettingZoo, \textbf{Ray RLlib}, MAPPO, Q-learning}

\entrydivider

\projectentry{Other Projects}{Computer Vision}
\begin{cvlist}
  \item \textbf{Marine Guardian} -- Geometric classifier + MobileNetV2 pipeline for ship detection in satellite imagery; 98.72\% classification accuracy, benchmarked against EfficientNet and ResNet-50 baselines.
  \item \textbf{Density-Based Crowd Counting} -- CSRNet trained in PyTorch for per-pixel density estimation in highly occluded surveillance scenes, avoiding failure modes of bounding-box detectors.
  \item \textbf{Facial Expression Benchmark} -- Controlled comparison of CNN, ResNet-34, and ViT backbones for emotion recognition; 87.5\% top-1 accuracy on the best configuration.
\end{cvlist}

% ════════════════════════════════════════════════════════════════════════════
%  PROFESSIONAL EXPERIENCE
% ════════════════════════════════════════════════════════════════════════════
\section*{Professional Experience}

\roleentry{Teaching Assistant, CSE 341: Computer Architecture}{January 2026 -- May 2026}{Department of Computer Science and Engineering, University at Buffalo}{Buffalo, NY}
\begin{cvlist}
  \item Conducted weekly tutorial sessions on \textbf{Assembly} and \textbf{VHDL} for 30+ students, live-coding solutions to architecture problems including cache design, memory hierarchies, and instruction pipelining.
  \item Led recitations walking through processor pipeline stages, hazard detection, and cache replacement policies; developed worked examples and problem sets to reinforce lecture material.
  \item Held weekly office hours; graded exams and assignments; supported 100+ students asynchronously on Piazza throughout the semester.
\end{cvlist}

\entrydivider

\roleentry{Tutor and Peer-Assisted Leader (PAL)}{August 2024 -- December 2025}{Tutoring and Academic Support Services (TASS), University at Buffalo}{Buffalo, NY}
\begin{cvlist}
  \item Served as \textbf{PAL for STA 119} (Introduction to Statistics): designed weekly worksheets aligned to recent lectures, then led two 50-minute collaborative sessions per week solving and explaining problems with students; course feedback indicated approximately 30\% improvement in quiz performance.
  \item Provided one-to-one 1-hour tutoring appointments across seven courses: \textbf{STA 119}, \textbf{CSE 115}, \textbf{CSE 116}, \textbf{CSE 220}, \textbf{MTH 141}, \textbf{MTH 142}, covering introductory programming, data structures, and calculus.
\end{cvlist}

\entrydivider

\roleentry{Founder, Inference Foundry}{May 2026 -- Present}{Open Research Software Initiative}{Remote}
\begin{cvlist}
  \item Launched an open-source initiative building reproducible ML inference tooling and shared contributor infrastructure for the scientific ML community (\href{https://inference-foundry.rweb.site}{inference-foundry.rweb.site}).
\end{cvlist}

\entrydivider

\roleentry{Founder and Vice President}{October 2023 -- May 2024}{UB National Student Data Corps (NSDC)}{Buffalo, NY}
\begin{cvlist}
  \item Co-founded the UB chapter; launched the chapter website, organized data science workshops, and built networking programs connecting undergraduates to research and industry opportunities.
\end{cvlist}

% ════════════════════════════════════════════════════════════════════════════
%  HONORS AND AWARDS
% ════════════════════════════════════════════════════════════════════════════
\section*{Honors and Awards}
\begin{cvlist}
  \item \textbf{PEARL Award} -- \$2,500 competitive undergraduate research grant, University at Buffalo Experiential Learning Network (November 2025)
  \item \textbf{Health Futures Challenge} -- Second place (\$1,000) for OralScan, University at Buffalo (Spring 2026)
  \item \textbf{Cybersecurity Excellence} -- Top 100 globally, Northeastern Cybersecurity C2C Finals; 10th in world finals (2025)
  \item \textbf{Collegiate Lockdown} -- Top two teams representing UB; 4th in finals (2025)
  \item \textbf{Dean's List} -- University at Buffalo (Fall 2023, Spring 2024, Fall 2024)
  \item \textbf{Presidential Scholarship} -- \$15,000 per year
\end{cvlist}

% ════════════════════════════════════════════════════════════════════════════
%  COMPETITIONS AND CHALLENGES
% ════════════════════════════════════════════════════════════════════════════
\section*{Competitions and Challenges}
\begin{cvlist}
  \item IAA AppXcelerate Application, University at Buffalo (January--March 2025) --  OralScan
  \item Russell L.\ Agrusa CSE Student Innovation Competition (November 2024) -- team member, OralScan
  \item Aging Innovations Challenge (November 2024) -- team member, OralScan
  \item Community Champions for Disability Health Challenge (October 2024) --  OralScan
\end{cvlist}

% ════════════════════════════════════════════════════════════════════════════
%  TECHNICAL SKILLS
% ════════════════════════════════════════════════════════════════════════════
\section*{Technical Skills}
\begin{cvlist}
  \item \textbf{Languages:} Python, JavaScript, Java, C++, Rust
  \item \textbf{Machine learning and AI:} PyTorch, TensorFlow, computer vision, deep learning, graph neural networks (GNNs), YOLOv8, Vision Transformer (ViT), reinforcement learning (MAPPO, Q-learning), Ray RLlib
  \item \textbf{Scientific ML and materials:} ALIGNN, machine learning interatomic potentials (MLIPs), MACE, E(3)-equivariant MPNNs, Weights \& Biases, PyTorch Geometric
  \item \textbf{Software engineering:} React, React Native, Node.js, Django, Git, CUDA, OpenCV, NumPy, Matplotlib
  \item \textbf{Databases and cloud:} MongoDB, SQL, Firebase, SQLite, AWS, Google Colab, Firebase Hosting
\end{cvlist}

% ════════════════════════════════════════════════════════════════════════════
%  RESEARCH INTERESTS
% ════════════════════════════════════════════════════════════════════════════
\section*{Research Interests}
\begin{cvlist}
  \item Spatiotemporal machine learning and deep learning-based weather forecasting (DLWF)
  \item AI adversarial robustness and out-of-distribution generalization
  \item Geometric and equivariant deep learning; graph neural networks for structured data
  \item Computer vision and sequence modeling for healthcare and scientific applications
  \item Machine learning for materials science and scientific discovery
\end{cvlist}

\end{document}